\documentclass[preview]{standalone}

\usepackage{amsmath}
\usepackage{amssymb}
\usepackage{stellar}
\usepackage{definitions}

\begin{document}

\id{psicologia-valutazione-intelligenza}
\genpage

\section{Valutazione dell'intelligenza}

\begin{snippet}{valutazione-intelligenza-storia}
    I primi strumenti di misura dell'intelligenza risalgono alla seconda metà
    dell'Ottocento. Galton cercò di dimostrare l'ereditarietà dell'intelligenza
    misurando caratteristiche fisiche e psicologiche come statura, acutezza
    visiva e tempi di reazione. Binet criticò questo approccio e propose prove
    più direttamente legate alle capacità intellettive.
\end{snippet}

\begin{snippet}{binet-simon}
    Il test di Binet-Simon valutava abilità come comprensione, immaginazione,
    memoria e attenzione. Era finalizzato a identificare bambini con difficoltà
    di apprendimento, così da poterli seguire con percorsi educativi adeguati.
\end{snippet}

\begin{snippetdefinition}{eta-mentale-definition}{Età mentale}
    L'\emph{età mentale}, o \(EM\), indica il livello di capacità cognitive di
    un individuo, equiparato al rendimento medio dei soggetti della stessa età
    cronologica.
\end{snippetdefinition}

\begin{snippetdefinition}{eta-cronologica-definition}{Età cronologica}
    L'\emph{età cronologica}, o \(EC\), è l'età anagrafica effettiva di un
    individuo.
\end{snippetdefinition}

\begin{snippetexample}{eta-mentale-cronologica-example}{Età mentale ed età cronologica}
    Se un bambino di \(4\) anni supera le prove previste per la sua classe di
    età, si ipotizza che la sua età mentale coincida con l'età cronologica. Se
    invece risponde correttamente solo alle prove dei \(3\) anni, si suppone che
    la sua età mentale sia inferiore all'età cronologica.
\end{snippetexample}

\begin{snippetdefinition}{quoziente-intelligenza-definition}{Quoziente di intelligenza}
    Il \emph{quoziente di intelligenza}, o \(QI\), è un indice del funzionamento
    intellettivo. Nella formulazione originaria di Stern era ottenuto dal
    rapporto tra età mentale ed età cronologica, moltiplicato per \(100\):
    \[
        QI = \frac{EM}{EC} \times 100
    \]
\end{snippetdefinition}

\begin{snippetdefinition}{qi-deviazione-definition}{QI-deviazione}
    Il \emph{QI-deviazione} è un indice calcolato in base alla distanza del
    punteggio di un individuo dalla media del proprio gruppo di età, espressa in
    unità di deviazione standard.
\end{snippetdefinition}

\begin{snippet}{distribuzione-qi}
    I punteggi di \(QI\) tendono a distribuirsi secondo la classica forma a
    campana, con la maggior parte dei punteggi vicina alla media e valori estremi
    progressivamente meno frequenti.
\end{snippet}

\includesnpt[width=55\%,src=/snippet/static-psychology/iq-normal-distribution.jpg]{centered-img}

\begin{snippetdefinition}{teoria-chc-definition}{Teoria CHC}
    La \emph{teoria Cattell-Horn-Carroll}, o \emph{CHC}, è una tassonomia delle
    abilità cognitive che integra la teoria dell'intelligenza fluida e
    cristallizzata con la teoria a tre strati di Carroll.
\end{snippetdefinition}

\begin{snippet}{strati-chc}
    La teoria CHC prevede tre strati: uno strato superiore corrispondente al
    fattore \(g\), uno strato intermedio composto da abilità ampie come
    intelligenza fluida, intelligenza cristallizzata, memoria e velocità di
    elaborazione, e uno strato inferiore formato da numerosi fattori specifici.
\end{snippet}

\includesnpt[width=58\%,src=/snippet/static-psychology/chc-cognitive-abilities-model.jpg]{centered-img}

\begin{snippetexample}{wisc-iv-example}{WISC-IV}
    La \emph{WISC-IV}, quarta versione della Wechsler Intelligence Scale for
    Children, è rivolta a soggetti tra \(6\) e \(16\) anni. Comprende subtest
    organizzati in quattro indici: comprensione verbale, ragionamento
    visuo-percettivo, memoria di lavoro e velocità di elaborazione. Questi
    indici concorrono alla definizione del quoziente intellettivo.
\end{snippetexample}

\includesnpt[width=62\%,src=/snippet/static-psychology/wisc-iv-cognitive-indexes.jpg]{centered-img}

\begin{snippetdefinition}{matrici-progressive-raven-definition}{Matrici Progressive di Raven}
    Le \emph{Matrici Progressive di Raven} sono prove non verbali di
    intelligenza basate su stimoli geometrici, considerate strumenti
    relativamente \emph{culture free} perché non richiedono acquisizioni
    scolastiche specifiche.
\end{snippetdefinition}

\begin{snippetexample}{raven-example}{Matrici Progressive di Raven}
    Nelle Matrici Progressive di Raven il soggetto osserva stimoli figurativi
    geometrici e deve scegliere, tra più alternative, lo stimolo che completa in
    modo pertinente il quesito proposto.
\end{snippetexample}

\includesnpt[width=45\%,src=/snippet/static-psychology/raven-progressive-matrices.jpg]{centered-img}
\includesnpt[width=45\%,src=/snippet/static-psychology/raven-nonverbal-matrix-item.jpg]{centered-img}

\begin{snippetdefinition}{affidabilita-qi-definition}{Affidabilità del QI}
    L'\emph{affidabilità del QI} è la misura in cui il punteggio di quoziente
    intellettivo è stabile e predittivo del funzionamento intellettivo dello
    stesso individuo a età diverse.
\end{snippetdefinition}

\begin{snippet}{stabilita-qi}
    L'affidabilità del \(QI\) aumenta con l'età. Nei bambini molto piccoli è più
    difficile ottenere valutazioni stabili; dopo i \(7\)-\(8\) anni la misura
    diventa più affidabile, anche grazie alla maggiore uniformità dei test.
\end{snippet}

\includesnpt[width=58\%,src=/snippet/static-psychology/iq-stability-over-time.jpg]{centered-img}

\end{document}
